\section{Delegating a bounded task}
\subsection{Risk and verifiability}

\begin{frame}{More autonomy can accumulate more errors}
  \begin{itemize}
    \item A long task requires several decisions to be correct in sequence.
    \item One local mistake can influence every later step.
    \item Breaking the work into parts makes intermediate results observable.
    \item Tests reduce risk only when they check something related to the real question.
  \end{itemize}

  \vfill
  Coding agents have improved quickly on long-task benchmarks. Our project still
  requires its own checks \citep{metr2025}.
\end{frame}

\begin{frame}{The frontier is irregular}
  Two tasks that look similar to us may fall on different sides of the system's
  capability.

  \begin{columns}[T,onlytextwidth]
    \column{0.48\textwidth}
      \begin{block}{Often easier to verify}
        Rename variables, change a format, add a stated assertion, or reproduce
        a documented transformation.
      \end{block}
    \column{0.48\textwidth}
      \begin{block}{Often harder to verify}
        Resolve an ambiguous policy concept, infer missing context, or decide
        whether an association justifies action.
      \end{block}
  \end{columns}

  \vfill
  Confidence gives us no information about where the task was on this ``jagged
  frontier'' \citep{dellacqua2026}.
\end{frame}

\begin{frame}{Match autonomy to risk and verifiability}
  \begin{columns}[T,onlytextwidth]
    \column{0.62\textwidth}
      \centering
      \begin{tikzpicture}[x=0.88cm,y=0.72cm,
        action/.style={align=center,text width=2.75cm,font=\small}]
        \draw[->,thick] (0,0) -- (7.2,0);
        \draw[->,thick] (0,0) -- (0,4.15);
        \draw[palegray] (3.55,0) -- (3.55,3.95);
        \draw[palegray] (0,1.95) -- (7.0,1.95);
        \node[anchor=north,font=\scriptsize] at (3.55,-0.42) {Result is easier to verify $\longrightarrow$};
        \node[rotate=90,anchor=south,font=\scriptsize] at (-0.50,1.95) {Harm if it fails};
        \node[action] at (1.78,2.92) {\color{crimson}\bfseries Person decides\\{\color{muted}\scriptsize policy judgment}};
        \node[action] at (5.30,2.92) {\bfseries Agent + gates\\{\color{muted}\scriptsize execute with review}};
        \node[action] at (1.78,0.94) {\bfseries Explore + review\\{\color{muted}\scriptsize questions and options}};
        \node[action] at (5.30,0.94) {\color{crimson}\bfseries Delegate + test\\{\color{muted}\scriptsize mechanical change}};
      \end{tikzpicture}
    \column{0.34\textwidth}
      \small
      \textbf{Verifiability} is high when a source, test, or independent
      comparison can expose the error.

      \vspace{0.45em}
      \textbf{Harm} includes cost, reversibility, and consequences for people,
      data, and claims.
  \end{columns}
\end{frame}

\subsection{A work order}

\begin{frame}{Write four things before execution}
  \begin{center}
  \begin{tikzpicture}[
    node distance=0.42cm,
    processbox/.style={draw=palegray,rounded corners=2pt,text width=2.8cm,
      minimum height=1.55cm,align=center,font=\small},
    arr/.style={-{Latex[length=2mm]},color=crimson,thick}
  ]
    \node[processbox] (goal) {\textbf{1. Goal}\\What concrete result do I want?};
    \node[processbox,right=of goal] (context) {\textbf{2. Context}\\Which files and definitions matter?};
    \node[processbox,right=of context] (permission) {\textbf{3. Permission}\\What may the agent read, change, or run?};
    \node[processbox,right=of permission] (proof) {\textbf{4. Proof}\\What evidence will show completion?};
    \draw[arr] (goal)--(context);
    \draw[arr] (context)--(permission);
    \draw[arr] (permission)--(proof);
  \end{tikzpicture}
  \end{center}

  \vfill
  Describe the proof before delegating the task.
\end{frame}

\begin{frame}{A weak request}
  \vspace{0.3em}
  \begin{block}{Request}
    Review my project and fix what is wrong.
  \end{block}

  What is missing?
  \begin{itemize}
    \item The object is unbounded.
    \item Relevant project context is missing.
    \item The permission boundary is missing.
    \item Evidence is undefined.
    \item Completion is undefined.
  \end{itemize}
\end{frame}

\begin{frame}[fragile]{A bounded, read-only work order}
\begin{lstlisting}
Goal: Explain how a student should start this project.
Context: Read START-HERE.md and the setup guide.
Permission: Read only. Do not edit files or browse the web.
Constraints: Separate facts in the files from your inferences.
Proof: List the steps in order, cite each source file, and end
with three checks that the student can run manually.
\end{lstlisting}

  \vfill
  The request is long because it makes the files, limits, and completion
  evidence explicit enough for the student to review.
\end{frame}

\begin{frame}{Useful context can change a decision}
  \begin{itemize}
    \item Give the agent the source file directly.
    \item Identify the current version and the source of truth.
    \item Separate facts, decisions, examples, and open questions.
    \item Ask for references to a file, line, table, or URL.
    \item Leave out material that cannot change the decision.
  \end{itemize}

  \vfill
  \begin{center}
    {\Large\color{crimson}Good context is information that can change a decision.}
  \end{center}
\end{frame}

\begin{frame}{Give permissions according to the task}
  \begin{columns}[T,onlytextwidth]
    \column{0.48\textwidth}
      \begin{block}{Reading}
        Web pages, documentation, code, public data, and metadata.
      \end{block}
      Reading is a good first step when we are still learning the project.
    \column{0.48\textwidth}
      \begin{block}{Action}
        Writing, running, deleting, publishing, sending, or buying.
      \end{block}
      Begin with less access. Increase it only when the task needs it.
  \end{columns}

  \vfill
  Never provide credentials, API secrets, identifiable student records, or
  confidential research material to an AI tool.
\end{frame}

\begin{frame}{Define evidence for the product}
  \small
  \begin{tabularx}{\textwidth}{@{}>{\bfseries}p{3.0cm}X@{}}
    \toprule
    Product & Minimum evidence \\
    \midrule
    Code & It runs from a clean start and relevant checks pass. \\
    Table or figure & Source, sample, units, and transformations agree. \\
    Explanation & It matches the files and distinguishes fact from inference. \\
    Text & Sources are attributed and no result is invented. \\
    \bottomrule
  \end{tabularx}

  \vfill
  A good check is cheap enough to repeat and strong enough to expose a plausible
  failure.
\end{frame}

\begin{frame}{Inspect the work behind the answer}
  \begin{center}
  \begin{tikzpicture}[
    node distance=0.26cm,
    processbox/.style={draw=palegray,rounded corners=2pt,text width=2.2cm,
      minimum height=1.25cm,align=center,font=\small},
    arr/.style={-{Latex[length=1.8mm]},color=crimson,thick}
  ]
    \node[processbox] (predict) {\textbf{Predict}\\What should happen?};
    \node[processbox,right=of predict] (delegate) {\textbf{Delegate}\\What may occur?};
    \node[processbox,right=of delegate] (observe) {\textbf{Observe}\\What did it touch?};
    \node[processbox,right=of observe] (verify) {\textbf{Verify}\\What check bears on it?};
    \node[processbox,right=of verify] (explain) {\textbf{Explain}\\Keep, revise, or reject};
    \draw[arr] (predict)--(delegate);
    \draw[arr] (delegate)--(observe);
    \draw[arr] (observe)--(verify);
    \draw[arr] (verify)--(explain);
  \end{tikzpicture}
  \end{center}

  \vfill
  ``Reject'' and ``still unresolved'' are valid outcomes. Keep a result when the
  available evidence supports it.
\end{frame}

\begin{frame}{Live read-only demonstration}
  Watch Codex read the course project without changing it.

  Record four things:
  \begin{enumerate}
    \item Which files did it read?
    \item Which actions did it request?
    \item Which statements came directly from files?
    \item Which decision still belongs to us?
  \end{enumerate}

  \vfill
  The demonstration uses the same project that you will open in Lab 1.
\end{frame}
